\FifthArticle{16}
As an example, take the equation
\[
x^7+28x^4-480=0
\]
for which
\[
\lambda=\frac{6750}{823543},\quad \log\frac{1}{\lambda}=2{,}0863825
\]
The equation has the first form and hence \textit{one} positive root; since $\lambda$ is smaller than $8$, it belongs to the second case. The first operation consists in satisfying the equation $\log\frac{1}{\lambda}=4A-3B$, or, if the calculation is begun with three decimal places, the equation
\[
2{,}086=4A-3B
\]
Even a quick glance at the table above shows that $A$ is to be sought between $0{,}5$ and $0{,}6$. For
\[
\begin{array}{c|c|c}
A&4A-3B&\mbox{Error}\\\hline
0{,}5&1{,}643&-0{,}443\\
0{,}6&2{,}109&+0{,}023
\end{array}
\]
which suggests the more accurate value $0{,}595$. A new calculation with the tables of five decimal places, thus putting $\log\frac{1}{\lambda}=2{,}08638$, gives
\[
\begin{array}{c|c|c}
A&4A-3B&\mbox{Error}\\\hline
0{,}595&2{,}08501&-0{,}00137\\
0{,}596&2{,}08961&+0{,}00323
\end{array}
\]
from which the still more accurate value $0{,}5953$ is found. Finally, for seven decimal places,
\[
\begin{array}{c|c|c}
A&4A-3B&\mbox{Error}\\\hline
0{,}5952&2{,}0859279&-0{,}0004546\\
0{,}5953&2{,}0863885&+0{,}0000060
\end{array}
\]
The value $A=0{,}5953$ therefore still requires the correction $-\frac{60}{4606}$ units in the fourth decimal place. I retain it in this form, since, if $x$ is to be determined by the first formula
\SourceLeaf{101}
\[
x^{7}=\frac{f}{c}
\]
all that is needed is to find the corresponding value of $C$. This is obtained by adding to the value listed beside $A=0{,}5953$, namely $C=0{,}6935705$, the correction $-\frac{60}{4606}\times798$, the latter being understood in units of the seventh decimal place. Thus
\[
\begin{array}{rcl}
C&=&0{,}6935695\\
\log f&=&2{,}6812412\\ \hline
7\log x&=&1{,}9876717\\
\log x&=&0{,}2839531\\
x&=&1{,}9228841
\end{array}
\]

To find the negative roots, put $x=-y$ and seek the positive roots of the equation
\[
y^{7}-28y^{4}+480=0
\]
This has the third form. Since $\frac{1}{\lambda}=\frac{823543}{6750}$ is greater than $\frac{7^{7}}{3^{3}4^{4}}=\frac{823543}{6912}$ but smaller than $2^{7}=128$, and at the same time $m$ is greater than $n$, the third case applies: there are two roots, to find which the equation
\[
2{,}0863825=3A+7B
\]
must be satisfied. From the closing remark of Article 14 it is known that one value of $A$ must be smaller and the other greater than $\log\frac{4}{3}=0{,}12494$. Bounds for the values of $A$ also follow at once from the table above with logarithms of three digits, which gives:
\[
\begin{array}{c|c|r}
A&3A+7B&\mbox{Error}\\ \hline
0{,}0&2{,}107&+0{,}021\\
0{,}1&2{,}078&-0{,}008\\
0{,}2&2{,}084&-0{,}002\\
0{,}3&2{,}132&+0{,}046
\end{array}
\]
If logarithms of four digits are to be used first for a more accurate determination, then for the first solution one initially has
\SourceLeaf{102}
\[
\begin{array}{r|r|r}
A&3A+7B&\mathrm{Error}\\\hline
0{,}05&2{,}0869&+0{,}0005\\
0{,}06&2{,}0847&-0{,}0017
\end{array}
\]
The tables of five digits then give
\[
\begin{array}{r|r|r}
0{,}052&2{,}08667&+0{,}00029\\
0{,}053&2{,}08638&0
\end{array}
\]
and finally those of seven digits
\[
\begin{array}{r|r|r}
0{,}0529&2{,}0863943&+0{,}0000118\\
0{,}0530&2{,}0863660&-0{,}0000165
\end{array}
\]
Hence
\[
\begin{array}{rcl}
A&=&0{,}0529417\\
\log f&=&2{,}6812412\\
7\log y&=&2{,}7341829\\
\log y&=&0{,}3905976\\
-y&=x=&-2{,}4580892
\end{array}
\]
For the second solution, the calculation, carried out in a similar way, is as follows:
\[
\begin{array}{r|r|r}
A&3A+7B&\mathrm{Error}\\\hline
0{,}19&2{,}0843&-0{,}0021\\
0{,}20&2{,}0868&+0{,}0004
\end{array}
\]
\[
\begin{array}{r|r|r}
0{,}197&2{,}08627&-0{,}00011\\
0{,}198&2{,}08654&+0{,}00016
\end{array}
\]
\[
\begin{array}{r|r|r}
0{,}1975&2{,}0863805&-0{,}0000020\\
0{,}1976&2{,}0864082&+0{,}0000257
\end{array}
\]
\[
\begin{array}{rcl}
A&=&0{,}1975072\\
\log f&=&2.6812412\\\hline
7\log y&=&2{,}8787484\\
\log y&=&0{,}4112498\\
x&=&-2{,}5778036
\end{array}
\]
\SourceLeaf{103}
The equation used here as our example was deliberately chosen so that two of its roots differ little. In such a case, as already observed in Article\ 11 above, the series are of little use because of their very slow convergence. At least a weak analogy can also be seen in the indirect solution: the successive approximations progress a little more slowly for the two negative roots, which are precisely the two that differ little, than for the positive root. There is, however, an essential difference: the very slow convergence of the series occurs for all the roots, whereas in the indirect procedure the slower approximation, noticeable even then only to a small degree, occurs solely for the two roots that differ little.

